% ============================================================================
%  HFT Tick-to-Trade FPGA Project — Issues Log
%  Compile with:  pdflatex issues_log.tex
% ============================================================================
\documentclass[11pt,a4paper]{article}

\usepackage[margin=1in]{geometry}
\usepackage{xcolor}
\usepackage{listings}
\usepackage{booktabs}
\usepackage{longtable}
\usepackage{enumitem}
\usepackage{titlesec}
\usepackage{fancyhdr}
\usepackage[hidelinks]{hyperref}

% ---- Colours -------------------------------------------------------------
\definecolor{codebg}{RGB}{245,245,245}
\definecolor{codekw}{RGB}{0,0,180}
\definecolor{codecm}{RGB}{0,128,0}
\definecolor{sev-high}{RGB}{200,0,0}
\definecolor{sev-med}{RGB}{200,120,0}
\definecolor{sev-low}{RGB}{0,120,0}

% ---- Code listing style --------------------------------------------------
\lstset{
  backgroundcolor=\color{codebg},
  basicstyle=\ttfamily\footnotesize,
  keywordstyle=\color{codekw}\bfseries,
  commentstyle=\color{codecm}\itshape,
  breaklines=true,
  frame=single,
  framesep=4pt,
  rulecolor=\color{gray!40},
  showstringspaces=false,
  columns=fullflexible,
  language=Verilog
}

% ---- Header / footer -----------------------------------------------------
\pagestyle{fancy}
\fancyhf{}
\lhead{HFT Tick-to-Trade FPGA}
\rhead{Issues Log}
\cfoot{\thepage}

% ---- Severity badge helper ----------------------------------------------
\newcommand{\sev}[1]{%
  \ifstrequal{#1}{HIGH}{\textcolor{sev-high}{\textbf{[HIGH]}}}{%
  \ifstrequal{#1}{MED}{\textcolor{sev-med}{\textbf{[MEDIUM]}}}{%
  \textcolor{sev-low}{\textbf{[LOW]}}}}}
\usepackage{etoolbox} % provides \ifstrequal

\titleformat{\section}{\Large\bfseries}{\thesection.}{0.5em}{}
\titleformat{\subsection}{\large\bfseries}{\thesubsection}{0.5em}{}

% ============================================================================
\begin{document}

\begin{titlepage}
\centering
\vspace*{3cm}
{\Huge\bfseries HFT Tick-to-Trade FPGA\par}
\vspace{0.5cm}
{\LARGE Issues \& Resolutions Log\par}
\vspace{2cm}
{\large Abhishek\par}
\vspace{0.3cm}
{\large Target: Xilinx Alveo U200 (\texttt{xcu200-fsgd2104-2-e})\par}
{\large Tooling: Vivado 2025.2 \quad Host: Ubuntu 24.04, 18\,GB RAM (VM)\par}
\vfill
{\large Last updated: 12 June 2026\par}
\end{titlepage}

\tableofcontents
\newpage

% ----------------------------------------------------------------------------
\section{Overview}
This document tracks the engineering issues encountered while bringing up the
tick-to-trade pipeline (\texttt{itch\_parser} $\rightarrow$
\texttt{order\_book\_m2} $\rightarrow$ \texttt{strategy\_imbalance} $\rightarrow$
\texttt{latency\_counter}) through Xilinx Vivado synthesis. Each entry records
the \textbf{symptom}, the \textbf{evidence}, the \textbf{root cause}, and the
\textbf{resolution / status}.

\begin{center}
\begin{longtable}{@{}clll@{}}
\toprule
\textbf{\#} & \textbf{Issue} & \textbf{Severity} & \textbf{Status} \\
\midrule
\endhead
1 & Synthesis exhausts RAM; order book not inferring BRAM & \sev{HIGH}   & \textbf{Fixed} (distributed RAM) \\
2 & Vivado native crash during synthesis (\texttt{hs\_err\_pid}) & \sev{HIGH}   & Resolved by \#1 \\
3 & XDC constraint errors in \texttt{top.xdc} (lines 31--32)    & \sev{MED}    & \textbf{Fixed} \\
4 & Assorted synthesis warnings (set/reset, unconnected ports) & \sev{LOW}    & Cosmetic / review \\
\bottomrule
\end{longtable}
\end{center}

% ----------------------------------------------------------------------------
\section{Issue \#1 — Synthesis exhausts RAM; order book stored as flip-flops instead of BRAM \sev{HIGH}}

\subsection{Symptom}
A batch synthesis run (\texttt{synth\_design -top tick\_to\_trade\_top -part
xcu200-fsgd2104-2-e}) ran the VM out of memory. GNOME System Monitor showed
\textbf{18.0\,GB of 18.4\,GB (97.7\%) used}, swap climbing past 2.3\,GB, and the
disk pinned at $\sim$560\,MB/s (swap thrashing, not useful work). The run was
manually interrupted.

\subsection{Evidence}
From \texttt{Vivado/synth\_batch.log}:
\begin{lstlisting}[language=]
WARNING: [Synth 8-11357] Potential Runtime issue for 3D-RAM or RAM from
         Record/Structs for RAM entries_reg with 33280 registers
...
2 Input 1 Bit Muxes := 2218
5 Input 1 Bit Muxes := 521
...
Finished Cross Boundary and Area Optimization : elapsed = 00:14:51  free physical = 5393
Finished Applying XDC Timing Constraints      : free physical = 1690
Start Timing Optimization
INFO: [Common 17-41] Interrupt caught. Command should exit soon.
wait_on_runs: ... free physical = 167          <- 167 MB RAM left
\end{lstlisting}
A single optimization phase ran for $\sim$15 minutes, and free physical RAM fell
from $\sim$10\,GB to \textbf{167\,MB} before the interrupt.

\subsection{Root cause}
The order book stores 256 orders, each a 130-bit packed struct
($256 \times 130 = 33{,}280$ bits). Vivado mapped this to \textbf{33,280
individual flip-flops} plus a large read-multiplexer tree (the 2{,}218 + 521
muxes above) instead of a Block RAM. Building and optimizing that fabric is what
caused the runtime and memory blow-up.

The array fails BRAM inference because the reset branch clears every element at
once, which Block RAM physically cannot do:
\begin{lstlisting}
// rtl/order_book_m2.sv : 107-108  (inside the asynchronous reset)
for (int i = 0; i < ORDER_DEPTH; i++)
    entries[i] <= '0;
\end{lstlisting}
Seeing a requirement to reset all contents simultaneously, the synthesizer
falls back from BRAM to flip-flops.

\subsection{Resolution / Status \textnormal{(\textbf{fixed} 12 June 2026)}}
Applied to \texttt{rtl/order\_book\_m2.sv}:
\begin{enumerate}[leftmargin=*]
  \item Removed the \texttt{valid} field from the \texttt{order\_entry\_t} struct
        and the whole-array async reset loop, so the payload array
        \texttt{entries[]} no longer needs clearing on reset.
  \item Added a separate resettable bit-vector
        \texttt{logic [ORDER\_DEPTH-1:0] entry\_valid;} (256 FFs) — a slot's
        contents are trusted only when its valid bit is set, so the payload
        never needs a reset.
  \item Kept reads \emph{combinational} (0-cycle) so no read latency is added to
        this latency-critical pipeline. The payload therefore infers
        \textbf{distributed RAM (LUTRAM)} rather than block RAM; true BRAM was
        declined because it would require synchronous reads and an FSM pipeline
        change (+1 clock per access). See trade-off note below.
\end{enumerate}
Effect: storage drops from 33{,}280 flip-flops + huge mux trees to LUTRAM
$+$ 256 valid FFs; the multi-minute optimization phase and the OOM disappear.

\textbf{Verification:} all 15 \texttt{tb\_order\_book\_m2} and 4
\texttt{tb\_tick\_to\_trade} cocotb tests pass under Verilator with cycle
behaviour unchanged. (On a clean Verilator build pass
\texttt{EXTRA\_ARGS="-Wno-fatal"} to get past the pre-existing PINMISSING
warning from Issue \#4.)

\textbf{Trade-off note:} for true block RAM (\texttt{RAMB36}), convert all reads
(the C/D/E lookup and the RESCAN stream) to synchronous and add read-wait
pipeline stages; this adds one clock of read latency per access. Revisit if LUT
pressure on the target device becomes a concern.

% ----------------------------------------------------------------------------
\section{Issue \#2 — Vivado native crash during synthesis \sev{HIGH}}

\subsection{Symptom}
An earlier (GUI) synthesis attempt aborted, leaving eight JVM/native crash dumps
in the project root (\texttt{hs\_err\_pid75146.log}, \texttt{75251},
\texttt{75357}, \texttt{75462}, \dots).

\subsection{Evidence}
\begin{lstlisting}[language=]
# An unexpected error has occurred (6) Aborted
...
libstdc++.so.6 ... _ZSt10unexpectedv
libxv_common.so ... hdi::tcltasks::detail::collection_wrapper_base::~...
libxv_synth.so  ... HARTSInfo::periodicalDetect()
\end{lstlisting}
At GUI launch \texttt{vivado.log} reported \texttt{Available Virtual : 503 MB} —
the machine was already memory-starved before the crash. Multiple helper
synthesis processes aborted together (one dump per PID).

\subsection{Root cause}
A native abort (signal 6) inside the synthesis worker, occurring under severe
memory pressure. Given the timing and the near-zero available virtual memory,
this is most plausibly the GUI-side manifestation of Issue \#1: the same order
book FF/mux explosion drove the process out of memory and Vivado aborted rather
than exiting cleanly.

\subsection{Resolution / Status}
Treat as resolved-by-\#1: once the order book infers BRAM the memory pressure
disappears. If native aborts persist afterward, capture a fresh
\texttt{hs\_err\_pid} and investigate the Vivado install separately. The stale
crash dumps can be deleted from the project root.

% ----------------------------------------------------------------------------
\section{Issue \#3 — XDC constraint errors in \texttt{top.xdc} \sev{MED}}

\subsection{Symptom}
Two CRITICAL WARNINGs every run while parsing the timing constraints, meaning
some constraints are silently dropped.

\subsection{Evidence}
\begin{lstlisting}[language=]
CRITICAL WARNING: [Designutils 20-1307] Command 'remove_from_collection' is
   not supported in the xdc constraint file. [rtl/top.xdc:31]
CRITICAL WARNING: [Common 17-1548] Command failed: can't read "in_ports":
   no such variable [rtl/top.xdc:32]
\end{lstlisting}
The offending lines:
\begin{lstlisting}[language=]
# rtl/top.xdc : 30-32
set in_ports [remove_from_collection [all_inputs] [get_ports clk]]
set_input_delay -clock clk $IO_BUDGET $in_ports
\end{lstlisting}

\subsection{Root cause}
\texttt{remove\_from\_collection} is a Tcl/XDC command not permitted inside a
pure \texttt{.xdc} constraint file. Line 31 therefore fails, \texttt{in\_ports}
is never set, and line 32 then fails with ``no such variable'' — so the input
delay on all non-clock ports is never applied.

\subsection{Resolution / Status \textnormal{(\textbf{fixed} 12 June 2026)}}
Replaced the \texttt{remove\_from\_collection} line in \texttt{rtl/top.xdc} with
an XDC-legal \texttt{get\_ports} filter that selects every input port except the
clock:
\begin{lstlisting}[language=]
set_input_delay -clock clk $IO_BUDGET \
    [get_ports * -filter {DIRECTION == IN && NAME != clk}]
\end{lstlisting}
This removes both CRITICAL WARNINGs and ensures the input-delay budget is
actually applied (previously it was silently dropped).

\textbf{Verification:} the file was sourced under \texttt{tclsh} with the Vivado
commands stubbed (and \texttt{remove\_from\_collection} stubbed to raise) — it
now parses cleanly, whereas the old version errored on the unset
\texttt{in\_ports} variable. The two CRITICAL WARNINGs should be gone on the
next synthesis run.

% ----------------------------------------------------------------------------
\section{Issue \#4 — Assorted synthesis warnings \sev{LOW}}

These do not block synthesis but are worth a review pass:

\begin{itemize}[leftmargin=*]
  \item \textbf{Set/reset same priority} — \texttt{[Synth 8-7137]} on several
        \texttt{itch\_parser} registers (\texttt{msg\_type\_reg},
        \texttt{timestamp\_reg}, \texttt{order\_ref\_reg}, \texttt{side\_reg},
        \texttt{shares\_reg}, \texttt{price\_reg}; lines 139--144). Vivado warns
        this may cause sim/synth mismatch; clean up the reset coding style.
  \item \textbf{Unconnected ports} — \texttt{[Synth 8-7129]}:
        \texttt{m\_axis\_tready} on \texttt{tick\_to\_trade\_top}, and the upper
        bits of \texttt{s\_axil\_wdata}/\texttt{s\_axil\_wstrb} on
        \texttt{latency\_counter} have no load.
  \item \textbf{Constant-driven outputs} — \texttt{[Synth 8-3917]}:
        \texttt{m\_axis\_tdata[71:65]} tied to constant 0 (upper trade-output
        bits unused).
\end{itemize}
Action: confirm each is intentional (it is normal for unused AXI/AXIS sidebands
during bring-up); otherwise wire up or trim the interface.

% ----------------------------------------------------------------------------
\section{Environment notes}
\begin{itemize}[leftmargin=*]
  \item Host: \texttt{abhishek-VMware-Virtual-Platform}, Ubuntu 24.04.2,
        AMD Ryzen 7 4800H, \textbf{18{,}412\,MB} RAM, 12{,}884\,MB swap.
  \item Tool: Vivado v2025.2 (64-bit), SW build 6299465.
  \item Target part: \texttt{xcu200-fsgd2104-2-e} (Alveo U200) — a large
        datacenter device; its part/timing database alone is memory-heavy to
        load, which compounds Issue \#1.
\end{itemize}

\end{document}
